\documentclass[11pt]{article}
\usepackage{amsmath,amssymb,amsthm,geometry,hyperref}
\geometry{margin=1in}
\newcommand{\per}{\operatorname{per}}
\newcommand{\tr}{\operatorname{tr}}
\newcommand{\ran}{\operatorname{ran}}
\newtheorem{theorem}{Theorem}
\newtheorem{lemma}[theorem]{Lemma}
\title{A sum-of-squares proof of the three-pairs case\\of Marcus's block-permanent inequality}
\author{Research proof for independent verification}
\date{6 November 2026}
\begin{document}
\maketitle
\begin{abstract}
We give an explicit positive-integer sum-of-squares identity for three complex
quadratic forms in at most five variables. A projection argument then proves
Marcus's block-permanent inequality for every $6\times6$ PSD Hermitian matrix
partitioned into three $2\times2$ blocks, with the expected equality condition.
The unrestricted problem, for arbitrary block count and block size, remains
unresolved in this work. Novelty has not been independently assessed, and
Lean formalization has not been attempted.
\end{abstract}

\begin{theorem}\label{main}
Let $A$ be a complex Hermitian positive semidefinite matrix of order six,
partitioned into $2\times2$ blocks $A_{ij}$, $1\leq i,j\leq3$. Set
$G_{ij}=\per(A_{ij})$. Then
\[
 \per(A)\geq\per(G).
\]
If $A$ has no zero row, equality holds if and only if $A_{ij}=0$ for $i\ne j$.
In particular this proves the requested equality clause when the three
diagonal blocks are positive definite.
\end{theorem}

\section{Notation and the quadratic identity}
The Fock inner product on complex polynomials is conjugate-linear in its first
argument and satisfies
\[
 \langle z^\alpha,z^\beta\rangle_F=\mathbf1_{\alpha=\beta}\alpha!.
\]
It is equivalently integration against independent circular complex Gaussians
of variance one. For complex symmetric $r\times r$ matrices $Q_1,Q_2,Q_3$, set
\[
 p_i(z)=z^{\mathsf T}Q_i z,\quad
 G_{ij}=2\tr(Q_i^*Q_j)=\langle p_i,p_j\rangle_F,\quad
 F=\|p_1p_2p_3\|_F^2-\per(G).
\]
All matrix and tensor norms below are the ordinary sums of squared absolute
values of their entries. A star denotes the Hermitian adjoint.

For each choice of a distinguished matrix $C$ among $Q_1,Q_2,Q_3$, let $A,B$
be the other two, in either order. Define
\begin{align*}
 T_{abcd}&=A_{ab}(BC^*)_{cd}+B_{ab}(AC^*)_{cd},\\
 S_{abcd}&=T_{abcd}+T_{acbd}+T_{bcad},\qquad H=3T-S,\\
 W&=AC^*B+BC^*A,\\
 L&=A\tr(BC^*)+B\tr(AC^*).
\end{align*}
The matrices $W,L$ are symmetric. For a symmetric matrix $X$, define
\begin{align*}
 E(X)_{abcd}&=\delta_{ad}X_{bc}+\delta_{bd}X_{ac}+\delta_{cd}X_{ab},\\
 K(X)_{abcd}&=\delta_{ad}X_{bc}+\delta_{bd}X_{ac}-2\delta_{cd}X_{ab}.
\end{align*}
The summation $\sum_C$ always means the three choices of distinguished matrix.

\begin{lemma}[Raw contraction identity]\label{raw}
In every dimension,
\begin{equation}\label{rawidentity}
 27F=\sum_C\left(112\|S\|^2+16\|H\|^2+288\|W\|^2
                     +216\operatorname{Re}\langle W,L\rangle\right).
\end{equation}
\end{lemma}
\begin{proof}
Expand the six Gaussian holomorphic indices against the six antiholomorphic
indices. A contraction is described by a $3\times3$ nonnegative integer table
$N$ with all row and column sums two: its $(i,j)$ entry counts contractions
from $Q_i$ to $\overline{Q_j}$. Symmetry of the matrices makes all contractions
with the same table identical. The multiplicity of such a table is
$64/\prod_{ij}N_{ij}!$.

The six tables $N=2P$ for permutation matrices $P$ are exactly the terms of
$\per(G)$; their total multiplicity is $48$. The remaining fifteen tables
have total multiplicity $672$. The following table records the expansion of
each expression, with coefficients already summed over the three choices of
$C$. Entries in the representative column list the nine entries of $N$ row
by row. The last column is the coefficient of
$\sum_C 2\operatorname{Re}\langle W,L\rangle$.
\[
\begin{array}{c|c|r|r|r|r|r}
\text{number}&N\text{ representative}&F&\sum\|S\|^2&\sum\|H\|^2&\sum\|W\|^2&\text{cross}\\\hline
6&002110110&32&3&6&0&4\\
1&011101110&64&0&0&6&0\\
3&011110101&64&12&-12&2&0\\
3&101020101&32&6&12&0&0\\
2&101110011&64&18&-18&0&0
\end{array}
\]
For every row, multiplying the last four columns by $112,16,288,108$
respectively gives $27$ times the $F$ column. No other tables occur.
This proves \eqref{rawidentity}. The table is also reproduced by the exact
symbolic contraction enumerator \texttt{marcus-quadratic-identity.mjs},
which checks each of the fifteen tables individually, with integer arithmetic.
\end{proof}

\begin{lemma}[Positive identity in dimension five]\label{sos}
For three symmetric matrices of order five,
\begin{align}\label{sosidentity}
189F=\sum_C\bigl(&16\|7S-E(L+2W)\|^2
             +7\|4H+K(L-W)\|^2\nonumber\\
             &+378\|W\|^2+126\|2L+5W\|^2\bigr).
\end{align}
Consequently $F\geq0$ for three quadratic forms in at most five variables.
If all three forms are nonzero, equality holds precisely when their matrix
ranges are pairwise orthogonal.
\end{lemma}
\begin{proof}
The tensor $S$ is symmetric in its first three indices. Its trace is
\[
 \sum_c S_{abcc}=L_{ab}+2W_{ab}.
\]
The tensor $H$ is symmetric in its first two indices, its cyclic sum over the
first three indices is zero, and
\[
 \sum_c H_{abcc}=2(L-W)_{ab}.
\]
Direct contraction of Kronecker deltas gives, in dimension $r$,
\begin{align*}
 \|E(X)\|^2&=3(r+2)\|X\|^2,&
 \langle S,E(X)\rangle&=3\langle L+2W,X\rangle,\\
 \|K(X)\|^2&=6(r-1)\|X\|^2,&
 \langle H,K(X)\rangle&=-6\langle L-W,X\rangle.
\end{align*}
For the last identity, the cyclic relation for $H$ says that the symmetric
part of either of its other traces is minus half of the displayed trace.
Taking $r=5$ therefore yields the exact formulas
\begin{align*}
 \|7S-E(L+2W)\|^2&=49\|S\|^2-21\|L+2W\|^2,\\
 \|4H+K(L-W)\|^2&=16\|H\|^2-24\|L-W\|^2.
\end{align*}
Substituting these into \eqref{sosidentity} and expanding the matrix squares
gives seven times \eqref{rawidentity}. This proves the identity and positivity.
Smaller-dimensional matrices can be padded with zeros to order five.

For equality, the positive coefficients force $W=0$ and then $L=0$ for
each distinguished matrix. The other squares force $S=H=0$, hence $T=0$.
Thus
\[
 A\otimes BC^*+B\otimes AC^*=0.
\]
If $A,B$ are linearly independent, this gives $BC^*=AC^*=0$. If $B=tA$
with $t\ne0$, it gives $2tA\otimes AC^*=0$, with the same conclusion.
Applying this for every $C$ shows $Q_iQ_j^*=0$ whenever $i\ne j$.
Since the matrices are symmetric, this is equivalent to orthogonality of
their ranges. Conversely orthogonal ranges allow disjoint orthonormal
coordinate sets for the three forms, and their Gaussian moments factor,
giving equality.
\end{proof}

\section{Two elementary multiplication facts}
For a polynomial $f$, let $f^*(\partial)$ denote the constant-coefficient
differential operator obtained by conjugating its coefficients and replacing
$z_j$ by $\partial_j$. The commutation rule
$[\partial_i,z_j]=\delta_{ij}$ gives the normal-ordering identity
\begin{equation}\label{normal}
 \|fg\|_F^2=\sum_\alpha\frac1{\alpha!}
       \| (\partial^\alpha f)^*(\partial)g\|_F^2.
\end{equation}
If $f$ is homogeneous of degree $d$, the terms $|\alpha|=d$ alone give
\begin{equation}\label{bombieri}
 \|fg\|_F^2\geq\|f\|_F^2\|g\|_F^2.
\end{equation}
If $g$ also has degree $d>0$, the $\alpha=0$ term additionally gives
\begin{equation}\label{two}
 \|fg\|_F^2\geq\|f\|_F^2\|g\|_F^2+|\langle f,g\rangle_F|^2.
\end{equation}
For quadratic forms $f=z^{\mathsf T}Q_1z$, $g=z^{\mathsf T}Q_2z$, the
remaining terms in \eqref{normal} give the exact gap
\begin{equation}\label{twogap}
 \|fg\|_F^2-\per[\langle p_i,p_j\rangle_F]_{i,j=1}^2
       =16\|Q_1^*Q_2\|^2.
\end{equation}
Thus equality in \eqref{two} for quadratics means orthogonal matrix ranges.

We also need the linear-factor identity
\begin{equation}\label{linear}
 \|fL_y\|_F^2-\|f\|_F^2\|y\|^2=\|D_y f\|_F^2,
 \quad L_y(z)=\sum_s y_s z_s,\quad D_y=\sum_s\overline{y_s}\partial_s.
\end{equation}
If $f$ is a nonzero product of linear forms whose coefficient vectors span
$U$, then $D_y f=0$ exactly when $y\perp U$. Indeed, under translation in
the direction defining $D_y$, the degree in the translation parameter of
that product equals the number of factors with nonzero directional derivative.
It can be zero only if every factor has zero directional derivative.

\section{Projection to at most four variables}
Choose Gram vectors $v_{i1},v_{i2}$ for the six rows of the original matrix,
and set $p_i=L_{v_{i1}}L_{v_{i2}}$. Gaussian moment expansion gives
\[
 \per(A)=\|p_1p_2p_3\|_F^2,\qquad G_{ij}=\langle p_i,p_j\rangle_F.
\]
Let $U$ be the span of the four coefficient vectors in the first two blocks.
It has dimension at most four. In orthogonal coordinates $z\in U$ and
$w\in U^\perp$, expand
\[
 p_3(z,w)=p_0(z)+\sum_{\alpha\ne0}h_\alpha(z)\frac{w^\alpha}{\sqrt{\alpha!}}.
\]
The polynomial $p_0$ is the product of the projections of the two linear
factors onto $U$. Write $f=p_1p_2$ and
$q_{12}=\per[\langle p_i,p_j\rangle_F]_{i,j=1}^2$.
Orthogonality of the $w$ monomials, and expansion of the $3\times3$
permanent in its last diagonal entry, give the exact decomposition
\begin{align}\label{reduction}
\|fp_3\|_F^2-\per(G)
={}&\left(\|fp_0\|_F^2-
       \per\operatorname{Gram}(p_1,p_2,p_0)\right)\nonumber\\
 &+\sum_{\alpha\ne0}
    \left(\|fh_\alpha\|_F^2-q_{12}\|h_\alpha\|_F^2\right).
\end{align}
The first bracket is nonnegative by Lemma~\ref{sos}, as all its forms use
at most four variables. For every remaining bracket, \eqref{bombieri}
and \eqref{two} give
\[
 \|fh_\alpha\|_F^2-q_{12}\|h_\alpha\|_F^2
 \geq (\|f\|_F^2-q_{12})\|h_\alpha\|_F^2\geq0.
\]
This proves the inequality of Theorem~\ref{main}.

\section{Equality in the six-vector problem}
Assume none of the six Gram vectors is zero. Thus all $p_i$ are nonzero.
For a product $L_vL_w$, its symmetric coefficient matrix is
$(vw^{\mathsf T}+wv^{\mathsf T})/2$; its range is exactly
$\operatorname{span}\{v,w\}$, including when $v,w$ are dependent.

If equality holds in \eqref{reduction}, its first bracket is zero.
If $p_0$ were nonzero, Lemma~\ref{sos} would force its matrix range to be
orthogonal to the ranges of both $p_1,p_2$, whose sum is $U$. But its range
lies in $U$, a contradiction. Hence $p_0=0$. Since $p_3\ne0$, at least one
$h_\alpha$ is nonzero. Equality in \eqref{reduction} then forces
$\|f\|_F^2=q_{12}$. Equation~\eqref{twogap} shows that the first two block
spans are orthogonal.

Write the coefficient vectors in the third block as $x+x'$ and $y+y'$,
with $x,y\in U$ and $x',y'\in U^\perp$. The equation $p_0=L_xL_y=0$
implies $x=0$ or $y=0$. Suppose $x=0$ and $y\ne0$. The first vector is
nonzero, so $x'\ne0$, and at least one coefficient $h_\alpha$ of outside
degree one is a nonzero scalar multiple of $L_y$.
Its bracket in \eqref{reduction} is a positive multiple of
$\|D_y f\|_F^2$ by \eqref{linear}. This cannot vanish: $y\in U\setminus\{0\}$
and the linear factors of $f$ span $U$. This is a contradiction. Interchanging
the two vectors handles the other case. Therefore $x=y=0$.

The third block span is thus orthogonal to $U$, and all three block spans are
pairwise orthogonal. This is exactly block diagonality of the Gram matrix.
Conversely block diagonality gives equality by factorization of the permanent.
This completes the proof of Theorem~\ref{main}.

\section{Verification record and scope}
The contraction identity was checked exactly on all fifteen contributing
incidence tables by \texttt{marcus-quadratic-identity.mjs}; its data are in
\texttt{marcus-quadratic-identity.json}. The dimension-five identity was also
checked by directly forming its integer-scaled tensors on forty independent
integer complex matrix samples. Those sample checks are supplementary; the
contraction count and delta identities above are the proof.

The initial coefficient discovery used \texttt{marcus-quadratic-sos.mjs} and
the installed Z3 solver. Its fitted coefficients are not used as a proof:
the final identity has the explicit positive integer coefficients in
\eqref{sosidentity}. No numerical optimization or unproved rank-reduction
claim is needed in the argument above.

This addresses the entire $m=3,k=2$ case. It does not establish Marcus's
conjecture for arbitrary $m,k$, so the original goal is still active.
The source status screen was
\href{https://yair-lavi.com/mincs-list/conjecture-9/index.html}{the current research page}
and \href{https://conjectures.io/problems/fulloldproblemscatalog-marcusblockpermanent-conjecture}{the chosen Conjectures.io entry}.
The general problem statement can also be found in
\href{https://arxiv.org/html/2202.01867v1}{Wanless's survey}.
\end{document}
